% ============================================================
% OP3_1a_bridge_recon_v1.tex  (PROPOSAL-ONLY RECONNAISSANCE)
% OP3.1a: the bare-P3 -> P4 bridge (= completing OP-grad).
% NO preprint edits, NO mechanism claims, NO fast-derivation
% promise. Exploration of the recorded P(X) route + alternatives.
% ============================================================
\section*{OP3.1a bridge reconnaissance (proposal-only)}

\subsection*{0. Scope honesty}
OP3.1a is explored as a programme: complete, test, and where
necessary extend the bridge from the bare action to the
gradient-ordered sector.
Nothing here derives the mechanism; the deliverables are
(i)~a precise map of what the article already records,
(ii)~one computed sharpening (a fluctuation-matching question),
(iii)~a draft admissible-extension-class definition, and
(iv)~clearly separated sub-tasks.

\subsection*{1. What the article already records (key finding)}
The $P(X)$ route is \emph{already in the text} as motivation, not
derivation: eq:PX\_action introduces
$S_E\supset\int d^4X\,P(\partial_A\Phi\,\partial_A\Phi)$ as the
minimal $O(4)$-invariant gradient-sensitive extension;
eq:PX\_minimum records the stationary-point condition
$P_X(u_*^2)=0$, $P_{XX}(u_*^2)>0$; eq:gradient\_branch and
eq:gradient\_vac record the resulting linear ordered background
$\Phi_0=u_AX^A+\Phi_c$, $u_Au_A=u_*^2$, with the $S^3$ of
degenerate ordered states and spontaneous $O(4)\to O(3)$.
The status paragraph there states explicitly: motivation for P4,
not a derived consequence of P1--P3; the specific $P$, the value
of $u_*$, and the relation $u_*\leftrightarrow(\mu,\lambda)$ are
OP-grad.
The formal OP-grad statement asks for the dynamical mechanism and
notes that the dimensional estimate (eq:u0\_dimensional) is
consistent.
Adjacent recorded structure: the broken-phase fluctuation tensor
$K^{AB}=\beta\,\delta^{AB}-\alpha\,n^An^B$ with the
Lorentzian-phase condition $\alpha>\beta>0$; the
``saddle-driven restructuring'' language for the $O(4)\to O(3)$
transition; admissibility properties S7/S9/P6.
Consequence: \textbf{OP3.1a is not ``invent a mechanism'' but
``complete and test the recorded $P(X)$ route (and its
competitors)''}.

\subsection*{2. Computed sharpening (proposal-level algebra;
to be written out and independently checked in v2)}
Expanding a pure $P(X)$ term to quadratic order around the linear
background ($\partial_A\bar\Phi=u_*n_A$), with
$\delta X=2u_*\,n\!\cdot\!\partial\chi+(\partial\chi)^2$:
\[
S^{(2)}\;\propto\;\int d^4X\Bigl[\,P_X(X_0)\,(\partial\chi)^2
+2P_{XX}(X_0)\,u_*^2\,(n\!\cdot\!\partial\chi)^2\Bigr],
\]
i.e.\ an effective fluctuation tensor
$\mathcal K^{AB}\propto P_X(X_0)\,\delta^{AB}
+2P_{XX}(X_0)\,u_*^2\,n^An^B$.
\emph{Observation (hedged).}
At the recorded stationary point ($P_X=0$, $P_{XX}>0$) this is
purely longitudinal and positive: the transverse sector has no
kinetic term from $P$ alone, and the longitudinal sign is
Euclidean-like.
By contrast, matching the form $\beta\delta^{AB}-\alpha n^An^B$
with $\alpha>\beta>0$ would require $P_X(X_0)>0$ and
$P_{XX}(X_0)<0$ at the background point --- at which the linear
profile is no longer a solution of the pure-$P$ equations.
\emph{Matching question M-($\alpha,\beta$) (the sharp sub-question
this recon proposes to record).}
Identify the additional structure that bridges the recorded
$P(X)$-stationary motivation and the recorded broken-phase tensor
$\alpha>\beta>0$: candidates include amplitude--potential coupling
(the $V(\Phi)$ sector along the profile), a mixed
$P(X,\Phi)$ dependence, the orientation-sector stiffness
($\mathcal C_n u^2(\partial n)^2$) entering at EFT level, or a
level distinction between the bare-fluctuation tensor and the
EFT tensor of the emergent-metric sector.
No inconsistency is asserted: the two tensors may legitimately
live at different levels; the task is to exhibit the bridge
explicitly.
A correct resolution must also be consistent with the recorded
saddle character of the ordering configuration (one
negative-signature gradient direction in the broken phase), which
already indicates that naive minimisation is not the selection
notion.

\subsection*{3. Sub-task structure for OP3.1a}
\emph{(a) Admissible extension class (definition task).}
Single real $\Phi$; locality; $O(4)$ and $\mathbb Z_2$;
first pass $P(X,\Phi)$ with standard reduction to
eq:ECT\_action in the small-$X$/weak-coupling regime;
Euclidean boundedness below ($X\ge0$ makes $P$-boundedness a clean
condition); optional second pass: $(\partial^2\Phi)^2$-type terms.
The class must be required to reproduce the recorded
ordered-branch EFT interface ($Z_u$, $\mathcal C_n$, $W(u)$) upon
expansion.
\emph{(b) Background existence within the class.}
Conditions on $P(X,\Phi)$ under which exact or controlled
ordered backgrounds exist (the pure-$P$ linear profile; its
deformations once $V$/$\Phi$-dependence is on).
\emph{(c) M-($\alpha,\beta$) resolution.}
As in \S2.
\emph{(d) Matching $u_*\leftrightarrow(\mu,\lambda)$.}
Only after (a)--(c); the recorded dimensional estimate is the
current placeholder.
\emph{(e) Selection.}
Hand-off to OP3.1c: dominance of the ordered sector in the
$S_0$-weighted measure; not part of OP3.1a proper.

\subsection*{4. Alternative/auxiliary classes (recorded for
completeness; all deferred)}
\emph{(E2) Modulated/Lifshitz-type ordering:} higher-gradient
terms with a negative-stiffness window select a modulation
wavevector and hence a direction; relation of stripe-like order to
the coarse-grained $(u,n_A)$ variables would need its own
analysis.
\emph{(E3) Boundary-data selection:} within the recorded
boundary-value formulation, ordered sectors selected by boundary
data; requires a principle for boundary data.
\emph{(E4) Measure/fluctuation selection:} ordering favoured at
one loop (order-from-the-measure); belongs with OP3.1c.
None is article-facing now.

\subsection*{5. What is NOT proposed}
No claim that the recorded $P(X)$ motivation is wrong; no claim
that any class derives the gradient condensate; no insertion; no
literature (candidate anchors for later page-level verification:
the ghost-condensate $P(X)$ literature and Lifshitz-point
literature); no status changes; no $u_*(\mu,\lambda)$ numbers.

\subsection*{6. Risk register}
Forbidden: ``ECT derives gradient condensation''; ``the $P(X)$
extension explains the Big Bang''; presenting M-($\alpha,\beta$)
as a contradiction in the article; minimisation language for the
ordering configuration.
Safe: ``recorded motivation; completion tasks separated'';
``matching question recorded''; ``saddle character consistent with
the recorded framework''.

\subsection*{7. Questions for GPT}
(Q1)~Approve the focus: v2 writes out the \S2 algebra in full
(candidate mini-lemma: fluctuation tensor of $P(X)$ around the
linear background, with the sign corollary), states
M-($\alpha,\beta$) precisely, and drafts the admissible-class
definition of \S3(a)?
(Q2)~Is M-($\alpha,\beta$), once verified, suitable for eventual
insertion as a recorded sharp sub-question of OP-grad/OP3.1a
(clearly marked as a question, not an inconsistency claim), or
proposal-only until resolved?
(Q3)~E2 (Lifshitz/modulated): compute a small comparison in v2 or
defer entirely?
(Q4)~Any objection to treating OP3.1a $\equiv$ OP-grad-completion
formally in future status echoes (they are already cross-referenced
in the implemented OP3.1 block)?
